\chapter{Natural-language proof sketch}
\label{chap:proof-sketch}

This chapter records the mathematical argument before assigning its obligations
to Lean modules. The detailed source audits are in
\href{https://github.com/tdwag123/PreissFormalization/blob/main/blueprint/reviews/rectifiability-audit.tex}{the rectifiability review}
and
\href{https://github.com/tdwag123/PreissFormalization/blob/main/blueprint/reviews/preiss-proof-audit.tex}{the uniform-measure review}.
They include the author's errata and proposed repairs of steps that need more
justification than the printed argument provides. None of these sketches is a
claim that the corresponding Lean proofs have been completed.

\section{The target and the normalization}
Fix natural dimensions $m,n$, and a locally finite Borel measure $\mu$ on
$\R^n$. Suppose the lower and upper $m$-densities agree and are positive and
finite at $\mu$-almost every point. We will prove that $\mu\ll\mathcal H^m$
and that $\mu$ is carried by countably many Lipschitz-parametrized pieces.
For Euclidean targets these pieces may equivalently be taken to be full
Lipschitz graphs, up to a Hausdorff-null remainder. The equivalence is a reusable
result; the difficult image-to-graph direction is not needed by the main proof.

We reuse FormalizingGMT's closed-ball, diameter-normalized density and its
arbitrary-scalar tangent predicate. For the analytic argument set
\[
 T_{x,r}(z)=(z-x)/r,\qquad D^m_{x,r}\mu=r^{-m}(T_{x,r})_\#\mu.
\]
At a positive finite density point, dimensional tangents and GMT tangents
are related by positive finite scalar multiplication. They are different sets.
Write $q_x=\lim_{r\downarrow0}\mu(B(x,r))/r^m$. When all tangents are uniform,
every dimensional tangent has the same ball constant $q_x$. Scaling $\mu$
by $q_x^{-1}$ permits the pointwise argument to use uniform measures with
unit-ball mass one. The Hausdorff normalization is handled by an explicit
positive constant on planes; no Euclidean-volume equality is assumed silently.

The cases $m=0$ and $m>n$ are separate. In dimension zero, continuity from
above on small balls identifies the density with $\mu\{x\}$. The measure is
therefore carried by its positive-mass atoms, a countable set because $\mu$
is sigma-finite. Such a measure is zero-dimensional rectifiable. If $m>n$,
finite upper density implies $\mu\ll\mathcal H^m$, while $\mathcal H^m$
vanishes on $\R^n$. Thus $\mu=0$. The rest assumes $0<m\leq n$.

\section{The rectifiability criterion needed at the end}
We first prove the Marstrand--Mattila implication independently of Preiss:
positive lower and finite upper density, together with flatness of every
tangent almost everywhere, imply rectifiability. This only assumes two-sided
density bounds, not an existing density or a common tangent plane.

Stratify into bounded Borel pieces on which the small-ball bounds are uniform.
Covering comparisons give finite Hausdorff measure and mutually absolutely
continuous restrictions of $\mu$ and $\mathcal H^m$ on each piece. Apply
Radon--Nikodym and differentiation on these finite restrictions, never on
ambient $\mathcal H^m$ when $m<n$. Transfer the tangent hypotheses to
$\nu=\mathcal H^m|F$ for each finite piece $F$.

At a point with two-sided density bounds, every positive sequence of radii
tending to zero yields dimensional blow-ups with a nonzero convergent subsequence. If every tangent is flat,
this forces weak linear approximation: at every sufficiently small radius
there is a plane near which almost all the mass lies, and every small ball
centred on that plane carries a definite positive mass. Failure of the first
clause contradicts compact-set portmanteau for a flat limit. Failure of the
second supplies moving centres converging to a point of the limiting plane;
a fixed smaller ball around that point contradicts open-set portmanteau.
Both clauses are indispensable.

It remains to show that finite-measure weakly linearly approximable sets are
rectifiable. Take the maximal rectifiable part by a countable maximizing
sequence. If a positive-measure purely unrectifiable part remains, weak
approximation passes to it: differentiation and finite upper density make the
discarded local mass $o(r^m)$.

We can uniformize weak approximation on compact subsets without assuming
measurability of arbitrary existential plane conditions. Replace strict upper
bounds by non-strict bounds and open lower-bound balls by closed balls, with
slack in the aperture. For fixed radius the resulting relation in point and
plane is closed, by lower semicontinuity of open-set masses and upper
semicontinuity of closed-ball masses. Projection along the compact space of
planes preserves closedness. Intersecting over radii gives closed good sets.
Shrinking and enlarging apertures relates them to the original definition.
Regularity then supplies the compact subsets used below.

Pure unrectifiability forces arbitrarily nearby pairs that escape any fixed
cone: otherwise local pieces would lie in Lipschitz graphs. Uniform lower
mass turns approximate tube concentration into actual tube containment of
the compact good set. A cone-escaping pair makes its approximating plane
have a direction with very small projection onto any prescribed comparison
plane. Linear algebra bounds the projected tube by a thin slab. A disjoint
almost-everywhere covering and the lower mass bounds show that the projection
onto each $m$-plane has zero Hausdorff measure.

Now use both directions of approximation on nested compact good sets
$G\subset F$. In a disk of a comparison plane, choose disjoint disks at scales
set by distance to the null projection of a compact piece
$G\cap\overline B(a,2\sigma)$; compactness gives nearest projected points.
Obtain the base tube bound at radius $2\sigma$, so it controls
$F\cap B(a,4\sigma)$, including the entire compact piece. With working
aperture $\beta=2\eta$, the hole radii satisfy $\rho_i\leq\beta\sigma$;
for $\beta<1/8$, the nearest witnesses lie in $B(a,\sigma/2)$.
This buffer avoids applying an open-ball estimate to boundary points.
Their radius powers have a
positive total comparable to the original scale to the power $m$. Discard
those cylinders that meet $F$ too deeply; their radius sum is controlled by
the small discarded mass. Each remaining empty cylinder, together with
plane-to-$F$ approximation at a nearest projected point, gives a long column
of disjoint balls centred on $F$. At the local scale
$r_i=\rho_i/(8\beta)=\rho_i/(16\eta)$ there are at least a positive constant times $1/\eta$
balls, each with its prescribed lower mass. The columns are mutually disjoint.
Choose the uniform lower-mass constant $\delta$ first, then
$\eta$ sufficiently small relative to it, then the finer approximation error
$\gamma$ and discarded-mass threshold $t$. Their total mass contradicts
the upper bound.
This local-scale formulation avoids an incorrect scale equality in the source.
The purely unrectifiable part is null, proving the set criterion and hence
the measure criterion by countable reassembly.

\section{Part A: obtaining one flat tangent}
FormalizingGMT already proves that every tangent at almost every positive
finite existing-density point is uniform. Density bounds also give tangent
existence. Starting with a uniform tangent in ambient dimension larger than
$m$, its support is not the whole space. A tangent at a nearest support point
to a hole lies in a half-space. The corrected barycentre argument produces
either the same measure or a further tangent supported in some hyperplane.
The hyperplane need not be the original half-space boundary.

Retain the tangent, normalization, and isometry witnesses through this
reduction. An existential lower-dimensional measure without those witnesses
is insufficient. The support-centred scaling orbit is closed under iterating
these limits, by support-point approximation, joint continuity of blow-ups,
and diagonal selection. Induction on ambient codimension ends with a
full-dimensional uniform measure, already classified in GMT as a positive
multiple of volume. Transport back to obtain a flat pseudotangent of the
original uniform measure.

To turn this into a tangent of $\mu$ at the original point, prove a single
almost-everywhere recentering assertion, with all tangents, support centres,
and positive radii quantified inside the almost-everywhere binder. A
countable compact-test metric permits a contradiction using a positive-mass
class of bad centres. Such classes need not be Borel. First localize them to finite mass, then take
their minimal measurable hulls: a measurable set disjoint from the original class has zero
intersection mass with its hull. At density-one points of the hull, this
forces support points of a blow-up to be approximated by original bad centres.
Joint continuity and the metric separation defining the class contradict one
another. This is an explicit repair of the measurability issue noted in the
errata. Recenter the pseudotangent sequence and use closure of the tangent
set. At least one flat tangent exists almost everywhere.

\section{Uniform measures at infinity}
We need three facts about unit uniform measures: uniqueness of their limit
at infinity, a gap between flat and nonflat conical limits, and rigidity when
the limit at infinity is flat. They are proved without using rectifiability.

The ball law identifies radial distributions and hence Gaussian integrals.
Polynomial growth gives integrable Gaussian-polynomial bounds and tails that
are uniform over the normalized family. These estimates extend vague
convergence to the unbounded tests needed for moments. Divide every Gaussian
moment by the Gaussian mass, so these moment forms are invariant under
positive scalar multiplication of the measure.

Translating a Gaussian by a support point gives an exponential identity.
A dominated Taylor expansion gives quantitative remainders at those points.
A separate finite-dimensional graded linear-algebra lemma upgrades this to
an expansion of every moment form: a weighted projection onto the space
spanned by support-polynomial evaluation vectors extends analytically across
zero. A finite basis of support evaluations controls its remainder in norm.
Use a finite coordinate-product model for the polynomial space; it avoids
incorrectly normalized symmetric-tensor bases. Uniqueness of Taylor
coefficients makes the expansions coherent across truncation orders.

The scaling identity for moments now identifies every Gaussian-polynomial
moment of a limit at infinity. Odd scaled moments vanish; even leading
coefficients carry the essential factorial from their Taylor expansion.
Two candidate limits have the same moments and hence are equal. For this
last step, first extend equality between nearby Gaussian rates by a dominated
series, then iterate. The algebra of Gaussian-weighted polynomials separates
points on the one-point compactification; weighted Stone--Weierstrass and
Gaussian integrability identify compact-test integrals. Ordinary polynomial
approximation on a compact set alone would not control the tails.
Compactness promotes subsequential uniqueness to full convergence. Dilation
invariance then makes the unique limit conical.

For a conical measure, scaling and the integer-power expansions force odd
moments to vanish. Its quadratic moment form is positive semidefinite, has
trace $m$, and agrees with squared norm on the support. Homogeneous radial
integration converts unit-ball distance energy to Gaussian distance energy.
A top-$m$ eigenspace minimizes this energy. Small energy forces every unit
point of that plane close to the support, by compactness and classification
of uniform measures supported on a plane. A support point sufficiently close
to the $m$th eigenvector forces that eigenvalue to be at least one. Trace $m$
then makes the form a projection and the measure flat. This proof works for
all positive $m$; the separate low-dimensional classification is unnecessary.

If the limit at infinity is flat on $V$, its Gaussian moments identify the
leading even coefficients of the original measure. The first support
identity gives a normal vector $b$ with
$|Q_Vx|^2=2\langle b,x\rangle$ on the support. Thus the support lies in a
bounded tube about $V$. A compactness argument also bounds holes along $V$.
Fourth-order support identities define an adjusted quadratic form and a
quadratic normal-valued map. Gaussian trace and fourth-moment identities,
with justified differentiation, Fubini, and limits, give a trace lower bound
by a positive multiple of $|b|^2$. The same trace is zero. Hence $b=0$,
the support lies in $V$, and the uniform measure is flat. The intermediate
trace identities are separate obligations in the revised graph.

\section{Part B and the conclusion}
At a point with density constant $q_x$, normalize as above so every dimensional
tangent is unit uniform. A compactly supported squared-distance functional,
minimized over the compact space of planes, is continuous under vague
convergence and along positive scaling. A flat tangent yields scales with
energy tending to zero. A hypothetical nonflat tangent has a nonflat conical
limit at infinity, by rigidity; the gap gives scales with energy bounded
away from zero.

Choose threshold-crossing scales between these sequences and pass to a
convergent subsequence of crossing measures. The limit has positive threshold
energy. The ratios of the adjacent flat scales to crossing scales must tend
to infinity: a bounded ratio would relate the limit to a flat measure by a
finite dilation. The choice of crossing then bounds the energy of every
large dilation of the crossing limit below the gap threshold. Its limit at
infinity is flat by the conical gap. Rigidity makes the crossing limit itself
flat, contradicting its positive energy. Thus every tangent is flat.

Intersect the almost-everywhere sets for density, uniformity, Part A, and
Part B. Apply the independently established Marstrand--Mattila measure
criterion, then include the zero and excess-dimensional cases. This is
Preiss' theorem. The stronger classical characterization with one common
plane almost everywhere is a separate reusable branch: graph area formula,
graph blow-up convergence, restriction/density transfer, and tangent
normalization prove its necessity direction without circular use of Preiss.
